% !TeX root = ../papers/nc-p3.tex
\section{一般単元剰余の二つの四次方程式と二次剰余の全点}
\label{sec:rank-fifteen-unit-quartic-towers}
\begin{theorem}\label{thm:rank-fifteen-unit-quartic-towers}
$k$ を標数3の代数閉体とし、$X=\bar x\ne0,1$ を固定する。
$c^2=X$, $b=1/c$, $d=b^3-b$ と置く。符号付き倍率の関数体の一表示は
\[
 D=v^3-v,\quad \bar y=X^3D^2,\quad
 R_{a,e}^2=f_{a,e}=(v-a-e b)\!\prod_{j\in\mathbb F_3\setminus\{a\}}(v-j)
 \qquad(a\in\mathbb F_3,\ e=\pm1)
\]
で与えられる。この被覆は $\bar y$ 上次数384、種数97であり、分岐値は $0,(1-X)^2,\infty$ である。
さらに $\eta_e=(\prod_a R_{a,e})/D$, $T=\eta_-+\eta_+$ とすると $k(T)$ は
$\bar y$ 上次数8の有理中間体である。その上の全被覆は次の二四次式の同時分解体である。
\[
 Z^4-2e b Z^2-8D\eta_e Z+b^2-8e bD,\qquad
 \eta_--\eta_+=2d/T,\qquad D=(T^2+4d^2/T^2)/4 .
\]
同時分解体の $k(T)$ 上の次数は48で、各四次式の分解体は次数12、Galois群は $A_4$ である。
\end{theorem}
\begin{proof}
$K=1/v$, $H=c(K-v)$ なら $U=K^2$, $V=H^2$ から二つの正則性座標が得られる。
$n\pm H/n$, $P_\pm$, $Q_\pm$ の六つの二乗に $v^2$ を掛けると、順に
$(1-v^2)(1\pm2cv)$、$2v(1+v)(1\pm c(1-v))$、$2v(1-v)(1\pm c(1+v))$ となる。
定数倍と符号の付け替えで上の六つのモニック三次式になり、定数は $k$ で開平できる。
初めの二根の和の半分から $n$ も復元できる。

$b\notin\mathbb F_3$ なので $\mathbb F_3$, $\mathbb F_3+b$, $\mathbb F_3-b$ は交わらない。
各三次式には他の五式が消えない単根 $a+e b$ があり、六平方類は独立である。
$v$ への64重被覆の分岐はこの九点と無限遠で、各慣性指数は2であるから
$2g-2=-128+10\cdot32=192$ となる。$v\mapsto\epsilon v+j$ は
$\bar y$ を保ち、$f_{a,e}$ を定数倍の置換に移すので、六変換とも持ち上がる。
三つの有限剰余類のうち元の素体は $D=0$、残る二類は $D=\pm d$ に写る。
$X^3 d^2=(1-X)^2$ であり、$v\mapsto D$ は有限点で非分岐である。

$f_{a,e}=D(1-e b/(v-a))$ より $\prod_a f_{a,e}=D^2(D-e d)$ である。
従って $\eta_e^2=D-e d$、$\eta_-^2-\eta_+^2=2d\ne0$ となり、
両根は $T$ の有理式で戻る。$D-d$, $D+d$ は異なる単根を持つので
$k(T)/k(D)$ は次数4であり、$\bar y$ への次数は8である。
$\sum_a f_{a,e}=e b$, $\sum_{a<a'} f_{a,e}f_{a',e}=2e bD$ は三次式の展開から従う。
三根 $R_{a,e}$ の符号を偶数個だけ反転した四つの和を取ると、
積 $D\eta_e$ が固定される。四つの和を根とするモニック四次式の展開が定理の式である。
四根のうち二つの和の半分から各 $R_{a,e}$ を戻せる。さらに
$f_{0,e}-f_{1,e}=-e b(2v-1)$ から $v$ が戻るので、同時分解体は全関数体である。
$k(T,v)/k(T)$ は、次数3の $k(v)/k(D)$ と次数4の $k(T)/k(D)$ が線型無関連なので次数3である。
この上では各三根の偶数符号変更が独立に作用し、合わせて $4\cdot4$ 個ある。
$v\mapsto v+1$ は二組の添字を同時に3巡回させ、各 $\eta_e$ を固定する。
従って全次数は48で、一方の四根に作用する像は $V_4\rtimes C_3=A_4$ である。
\end{proof}

\begin{corollary}\label{cor:rank-fifteen-unit-elliptic-quotient}
一方の四次式の根 $Z$ を選び、$\eta=\eta_e$ とする。点を一つ指定した関数体
$k(\eta,Z)$ は種数2であり、両符号とも同じ楕円曲線
\[
                 Y^2=(l+1)(l+1+Xl^2)
\]
の二重被覆になる。具体的には $s=e b$, $a=Z^2/s$,
$l=(Z^2+\eta Z)/s$, $Y=2a-l-1$ である。
\end{corollary}
\begin{proof}
標数3では四次式は $(Z^2-s)^2+D(\eta Z+s)=0$、
$D=\eta^2+s^3-s$ である。$A=Z^2$, $B=\eta Z$ を入れて $A$ 倍すると
$A(A-s)^2+(B^2+(s^3-s)A)(B+s)=0$ となる。
$A=sa$, $B=s(l-a)$ と変換すると $a^2-(l+1)a=Xl^2(l+1)$ に等しく、
平方完成して上の式を得る。二次因子の判別式は $1-X\ne0$、
$l=-1$ での値は $X\ne0$ なので三次式は平方自由である。
$Z^2=sa$ を添加すると $\eta=s(l-a)/Z$ も戻る。
楕円曲線上で $a=0$ は $(l,a)=(0,0),(-1,0)$ に限られる。
前者では $a-(l+1)$ が単元で $a$ の零位数は2である。
後者では $Y$ が一様化元で $l+1$ の位数が2なので、$a=(Y+l+1)/2$ は単純零である。
無限遠では $l,Y$ の極の位数は2,3で、$a$ は三位の極を持つ。
二重被覆の分岐点はこの単純零と無限遠の二点であり、種数は2となる。
さらに $t=l+1$, $W=Y/c$ で $W^2=t^3+(1+1/X)t^2+t$ となる。
標数3のHasse不変量は $t^2$ の係数なので、この楕円曲線が超特異となるのは
$X=-1$ のときに限る。$j=(X+1)^6/(X^4(1-X))$ である。
特に $X=-1$ で $l=i t-1$ と置くと $Y^2=-i(t-t^3)$ となり、
定数倍で既出の $W^2=t-t^3$ の関数体をそのまま利用できる。
\end{proof}

\begin{proposition}\label{prop:rank-fifteen-genus-thirteen-window}
上の六根を含む体を $E$ とする。四窓テンソルは $k(v)$ 上次数16、種数13の体
\[
 F=k(v)\big(\sqrt{(v-a)^2-b^2}\ (a\in\mathbb F_3),\ \eta_+\big)
\]
に係数を持つ表示を持ち、その表示でもKlein次数に関して斉次である。
\end{proposition}
\begin{proof}
同じ $a$ の両根を同時に反転する操作を対反転と呼ぶ。
$S$ は $1,-1$ の二対、$P$ は $0,-1$ の二対を反転し、位数4の群を生成する。
$R_{a,+}R_{a,-}$ は $D/(v-a)$ と上の平方根の積である。
三つの対積の符号指標と、三つの正号根の積の符号指標は独立である。
従ってこれらで生成した体の次数は16で、これが固定体である。
三つの二次式の単根は $a\pm b$ の六点で、$\eta_+^2=D-d$ の三単根もそのうちにある。
さらに無限遠で後者が奇数位数の極を持つ。他に分岐はなく各慣性次数は2なので、
$2g-2=-32+7\cdot8=24$ である。

添字 $t\in\{0,1,2,3\}$ の上下ビットを $t_H,t_L$ と書く。
正規形の倍率 $\gamma_{h,j}$ で $S$ の作用は $(-1)^{h_H}$ 倍、
$P$ の作用は $(-1)^{j_L}$ 倍である。これは $(L,M),(A,B),(C,D)$ が順に
$0,1,-1$ の各対に線型であることから従う。
四つの位置すべてで低ビットをWalsh基底 $e_{2r}+(-1)^q e_{2r+1}$ に変える。
第1、第3位置ではさらに添字 $t=2r+q$ の係数に
$\phi_t=R_{1,+}^{r}R_{0,+}^{q}$ を掛ける。この変換は関数体上可逆である。
$\beta=(0,3,2,1)$ と置くと、変換後の係数は
$a\mathbin\oplus c=b\mathbin\oplus d=h$ のとき
\[
 4\phi_a\phi_c\sum_{k=0}^1\ \,\sum_{j_L=h_L}
 (-1)^{(a_L+b_L)k+j_Hc_H+\beta(j)_Hd_H}\,
 \gamma_{2h_H+k,j}
\]
であり、それ以外は零となる。実際、旧低ビット $c_L,d_L$ で符号を足すと、
$\beta(j)_L=j_L$ により、二つの直交条件はいずれも新しい低ビット差を $j_L$ と等しくする。
この式での $S,P$ の符号を $\phi_a\phi_c$ が打ち消すので係数は $F$ に入る。
添字付きの基底変換なので、この表示に四次中心も移せば各始点の商の次元が保たれる。
\end{proof}

\begin{theorem}\label{thm:rank-fifteen-quadratic-residue-pairs}
数体 $L$ が3上に剰余次数2以下の素点を持つとする。任意の $x\in L$, $y\in L^\times$ で、
通常な階数15窓は標準Koszul左右AS正則で、各始点が四中心四次式上自由階数256である。
特に、$x,y$ が任意の有限個の二次数体の合成体に属すればよい。
\end{theorem}
\begin{proof}
素点の付値を窓の係数体に延長する。負付値の $x$ は前節、正付値または零の $x$ は
第\ref{sec:rank-fifteen-nonunit-pole-strip}節で決まる。単元の $x$ で $y$ が非単元なら
既知の正付値領域または極の定理による。単元同士で $\bar x=\pm1$ の場合も既に認証済みである。

残る $(X,Y)$ は $X\in\mathbb F_9\setminus\mathbb F_3$, $Y\in\mathbb F_9^\times$ であり、
同時Frobeniusで24軌道である。\path{computations/verify_rank_fifteen_unit_quartic_towers.py} は
各軌道で $c^2=X$, $(v^3-v)^2=Y/X^3$ を解き、倍率の方程式と二座標を検査する。
平方根の後のArtin--Schreier方程式は次数3以下で分解するため、$c,v$ は
$\mathbb F_{3^{12}}$ に取れ、残る平方根は $\mathbb F_{3^{24}}$ に取れる。
各表示で未切断経路を五次まで、文字別四次中心を計算し、四本の中心式を課して
全四始点で十三次まで標準商を確認する。単元座標でのモニックな持上げと
全次数の比較は定理\ref{thm:rank-fifteen-unit-residues}の証明と同じである。
二次数体の合成体では分解群の元の位数は高々2であり、その巡回な剰余体Galois群も位数2以下となる。
\end{proof}
